\chapter{Lessons Learned}\label{ch:lessons}


\section{Reflexion der Erfahrungen aus der Analyse}\label{sec:reflexion-erfahrungen}

Die Datenaufbereitung erwies sich als der eigentliche Engpass, selbst bei synthetischen Daten. Der Erfahrungswert, dass die Datenakquisition 60 bis 80 Prozent der Projektzeit beansprucht, bestätigte sich konkret \parencite[Folie~414]{coletta2026_3006}. In einer ersten Fassung des Simulationsmodells erzeugten nicht streng monotone Zeitstempel Schein-Kanten im entdeckten Modell. Solche Granularitäts- und Sortierprobleme treten auch bei realen \gls{sap}-Extraktionen auf, wenn tagesgenaue Felder wie \gls{augdt} mit sekundengenauen Zeitstempeln gemischt werden, und müssen vor der Discovery behoben werden, sonst analysiert man Artefakte statt Prozesse. Ebenso zeigte sich, dass die Methodenwahl kein akademisches Detail ist. Zwischen Alpha-Algorithmus und Inductive Miner lagen auf demselben Log deutliche Unterschiede (Fitness 0,646 gegenüber 0,996), und beim Soll-Abgleich machte erst der Wechsel vom Token-Replay zu Alignments die tatsächliche Konformität sichtbar (Kapitel~\ref{sec:gaps-auswirkungen}). Wer Werkzeugausgaben interpretiert, muss wissen, welcher Algorithmus mit welchen Parametern dahinterliegt. Eine zweite Grenze betrifft die Reichweite der Daten selbst. Ein Event-Log zeigt, \emph{was} geschehen ist, nicht \emph{warum}. In der Fallstudie von Adolfsson/Saleh blieb eine durchgängige manuelle Auftragsabstimmung mit dem Vertrieb unsichtbar, weil sie keine digitale Spur im \gls{erp}-System hinterließ \parencite[S.~50 und 62]{adolfsson2026o2c}. Für den synthetischen Log ist diese Grenze konstruktionsbedingt gegenstandslos, für die Übertragung auf reale Systeme jedoch zentral. Konkret zeigte sich die Bedeutung der Datenaufbereitung daran, dass erst die Erzwingung streng monoton steigender Zeitstempel je Fall die scheinbaren Rücksprünge im Modell beseitigte; ein Detail, das bei einer flüchtigen Betrachtung der Werkzeugausgabe leicht als echter Prozessschritt fehlgedeutet worden wäre. Die Lehre daraus ist, dass die Plausibilität des Event-Logs vor jeder inhaltlichen Interpretation zu prüfen ist.


\section{Erkenntnisse und Empfehlungen für ähnliche Projekte}\label{sec:erkenntnisse-empfehlungen}

Daraus folgen vier Empfehlungen: Erstens sollten Kennzahlen vor der Analyse definiert werden (Sperrquoten, Liegezeiten je Übergang, Happy-Path-Anteil, Zeit bis zur Faktura). Ohne Ziel-\gls{kpi}s erzeugt Process Mining beeindruckende, aber folgenlose Visualisierungen. Zweitens verdient die Log-Qualitätsprüfung (Monotonie, Granularität, Vollständigkeit) einen festen Platz vor jeder Discovery. Drittens gehört beim Übergang auf reale Logs der Datenschutz an den Anfang, mit Pseudonymisierung der Benutzerkennungen und früher Einbindung von Betriebsrat und Datenschutzbeauftragten (Kapitel~\ref{sec:risikobewertung}). Viertens gilt der Customizing-Leitsatz „so viel Standard wie möglich“ auch für die Behebung der Gaps. Alle sechs Maßnahmen kommen ohne Modifikation aus, was Wartbarkeit und Release-Fähigkeit erhält \parencite[Folie~248]{coletta2026_0505}. Diese Empfehlungen decken sich mit der Praxisliteratur, die rät, Analysen breit-aggregiert zu beginnen, das Werkzeug mit Interviews und Workshops zu kombinieren und den Untersuchungsumfang zunächst eng zu fassen \parencite[S.~71 f.]{adolfsson2026o2c}. Eine fünfte Erkenntnis betrifft das Zusammenspiel von Methode und Fachwissen: Process Mining liefert präzise Diagnosen, aber keine Therapien. Die Entscheidung etwa, dass die Batch-Fakturierung trotz ihrer scheinbaren Harmlosigkeit die lohnendste erste Maßnahme ist, erfordert betriebswirtschaftliches Urteilsvermögen und Kenntnis des \gls{sap}-Customizings. Das datengetriebene Vorgehen ersetzt die Prozessworkshops also nicht, es macht sie gezielter und beweisbarer \parencite[Folie~247]{coletta2026_0505}. Nicht zuletzt lohnt der Blick über den eigenen Prozess hinaus: Die Praxisliteratur und die herangezogene Fallstudie zeigen übereinstimmend, dass die größten Effekte weniger in ausgefeilteren Algorithmen als in sauberen Stammdaten und klar geschnittenen Prozessen liegen \parencites[S.~110 ff.]{nadukuru2024o2c}[S.~71 f.]{adolfsson2026o2c}. Process Mining ist damit vor allem ein Instrument, um diese unspektakulären, aber wirksamen Hebel überhaupt sichtbar und ihre Behebung messbar zu machen.


\section{Langfristige Auswirkungen auf das Unternehmen}\label{sec:langfristige-auswirkungen}

Langfristig liegt der Wert weniger in der einmaligen Diagnose als in der Verstetigung. Da Simulation und Analyse vollständig skriptbasiert und reproduzierbar sind, lässt sich dieselbe Auswertung nach jeder Prozessänderung erneut ausführen, als kontinuierliches Monitoring mit den Kennzahlen aus Anhang A3 (Kapitel~\ref{sec:modifikationen-monitoring}). Für die Organisation verschiebt das die Diskussionsgrundlage dauerhaft von Einzelmeinungen zu gemessener Evidenz. Prozessverbesserung wird überprüfbar und Prozessverantwortung messbar. Erfahrungsgemäß entsteht daraus ein kultureller Wandel: Wo Kennzahlen sichtbar sind, kehren alte Muster wie schleichend zunehmende manuelle Preisänderungen seltener unbemerkt zurück, weil jede Abweichung im Monitoring auffällt. Perspektivisch ist die Übertragung auf weitere Belegflüsse wie den Purchase-to-Pay-Prozess angelegt, für den der \gls{bpi}-Challenge-2019-Log als öffentliches Übungsfeld bereitsteht \parencite{vandongen2019bpi}.
